% Exam 3
\question\MatchQuestion{}{}
\question\MatchQuestion{CaM1}{Toolkit gene associated with variation in beak length in birds.}
\question\MatchQuestion{Co-option}{Existing gene evolves new function without duplication.}
\question\MatchQuestion{BMP / SHH}{Two toolkit genes associated with the development of scales and feathers in alligators and birds. \textbf{Write both for full credit.}}
\question\MatchQuestion{Serial homology}{Repeating segments within an organism, serving different functions, inherited from a common ancestor.}
\question\MatchQuestion{FOXP2}{Toolkit gene associated with vocal ability in humans and some birds.}
species.}
\question\MatchQuestion{Anisogamy}{Sexes have different sized gametes.}
\question\MatchQuestion{Muller's Ratchet}{Accumulation of detrimental mutations in asexually-reproducing species.}
\question\MatchQuestion{Allometry}{Different body parts develop at different rates.}
\question\MatchQuestion{Paedormophosis}{Retention of juvenile or larval characters in adults.}
\question\MatchQuestion{Maternal effect}{mRNA transcripts inside unfertilized egg.}
\question\MatchQuestion{End Permian}{Largest mass extinction.}
\question\MatchQuestion{End Cretaceous}{Most famous mass extinction.}
\question\MatchQuestion{End Ordovician}{First mass extinction.}



Persimmion trees produce either male flowers (no carpels) or female flowers (no stamens). Use the floral quartet model to explain which (if any) class or classes of MADS genes could be lost (no expression) to produce 1) male flowers and 2) female flowers. If the model does not work for one or both sexes of flowers, explain why for each sex.
